\documentclass[acmsmall,screen]{acmart}
%%
%% \BibTeX command to typeset BibTeX logo in the docs
\AtBeginDocument{%
  \providecommand\BibTeX{{%
    Bib\TeX}}}

%% Rights management information.  This information is sent to you
%% when you complete the rights form.  These commands have SAMPLE
%% values in them; it is your responsibility as an author to replace
%% the commands and values with those provided to you when you
%% complete the rights form.
\setcopyright{acmlicensed}
\copyrightyear{2018}
\acmYear{2018}
\acmDOI{XXXXXXX.XXXXXXX}

%%
%% These commands are for a JOURNAL article.
\acmJournal{JACM}
\acmVolume{0}
\acmNumber{0}
\acmArticle{0}
\acmMonth{0}

%%
%% Submission ID.
%% Use this when submitting an article to a sponsored event. You'll
%% receive a unique submission ID from the organizers
%% of the event, and this ID should be used as the parameter to this command.
\acmSubmissionID{XXX-XXX-XXX}

%%
%% For managing citations, it is recommended to use bibliography
%% files in BibTeX format.
%%
%% You can then either use BibTeX with the ACM-Reference-Format style,
%% or BibLaTeX with the acmnumeric or acmauthoryear sytles, that include
%% support for advanced citation of software artefact from the
%% biblatex-software package, also separately available on CTAN.
%%
%% Look at the sample-*-biblatex.tex files for templates showcasing
%% the biblatex styles.
%%

%%
%% The majority of ACM publications use numbered citations and
%% references.  The command \citestyle{authoryear} switches to the
%% "author year" style.
%%
%% If you are preparing content for an event
%% sponsored by ACM SIGGRAPH, you must use the "author year" style of
%% citations and references.
%% Uncommenting
%% the next command will enable that style.
%%\citestyle{acmauthoryear}


%%
%% end of the preamble, start of the body of the document source.
\begin{document}

%%
%% The "title" command has an optional parameter,
%% allowing the author to define a "short title" to be used in page headers.
\title{Public interest data science}

%%
%% The "author" command and its associated commands are used to define
%% the authors and their affiliations.
%% Of note is the shared affiliation of the first two authors, and the
%% "authornote" and "authornotemark" commands
%% used to denote shared contribution to the research.
\author{Brian C. Keegan}
\email{brian.keegan@colorado.edu}
\orcid{0000-0002-7793-398X}
\affiliation{%
  \institution{University of Colorado Boulder}
  \city{Boulder}
  \state{Colorado}
  \country{USA}
}

%%
%% By default, the full list of authors will be used in the page
%% headers. Often, this list is too long, and will overlap
%% other information printed in the page headers. This command allows
%% the author to define a more concise list
%% of authors' names for this purpose.
\renewcommand{\shortauthors}{Keegan}

%%
%% The abstract is a short summary of the work to be presented in the
%% article.
\begin{abstract}
Data-driven systems increasingly govern critical aspects of public life, yet access to the data and infrastructures needed to study and oversee these systems is rapidly shrinking. This article defines public interest data science as an emerging field that adapts the public missions of journalism, law, accounting, planning, and engineering to contemporary data-intensive environments. It argues that public interest data science arises in response to three structural dynamics: enclosure, where once-public data is privatized or locked down; exemption, where platforms and AI developers operate outside meaningful regulatory scrutiny; and erosion, where fragmentation and infrastructural decay undermine long-term observation and reproducible research. Against these forces, the article advances a constructive agenda organized around openness, oversight, and ownership as countervailing pillars of accountability and stewardship. Public interest data science is distinguished from adjacent efforts such as data for good, civic technology, critical data studies, and data justice by its emphasis on building durable infrastructures and regulatory architectures for democratic oversight. Three case studies on mining municipal archives, infrastructuring worker observatories, and designing Anthropocene data trusts illustrate how public interest data science reimagines the relationships between data, power, and the public good in practice. The article concludes by outlining how public interest data science can provide the foundations for future institutions analogous to current regulatory regimes.
\end{abstract}

%%
%% The code below is generated by the tool at http://dl.acm.org/ccs.cfm.
%% Please copy and paste the code instead of the example below.
%%
\begin{CCSXML}
<ccs2012>
   <concept>
       <concept_id>10003456.10003457.10003567.10010990</concept_id>
       <concept_desc>Social and professional topics~Socio-technical systems</concept_desc>
       <concept_significance>300</concept_significance>
       </concept>
   <concept>
       <concept_id>10003456.10003462.10003480.10003484</concept_id>
       <concept_desc>Social and professional topics~Technology and censorship</concept_desc>
       <concept_significance>300</concept_significance>
       </concept>
   <concept>
       <concept_id>10003456.10003462.10003544.10003589</concept_id>
       <concept_desc>Social and professional topics~Governmental regulations</concept_desc>
       <concept_significance>300</concept_significance>
       </concept>
   <concept>
       <concept_id>10010405.10010455</concept_id>
       <concept_desc>Applied computing~Law, social and behavioral sciences</concept_desc>
       <concept_significance>500</concept_significance>
       </concept>
   <concept>
       <concept_id>10010405.10010476.10010936</concept_id>
       <concept_desc>Applied computing~Computing in government</concept_desc>
       <concept_significance>300</concept_significance>
       </concept>
    <concept>
        <concept_id>10003120.10003130.10003131</concept_id>
        <concept_desc>Human-centered computing~Collaborative and social computing~Collaborative and social computing theory, concepts and paradigms</concept_desc>
        <concept_significance>300</concept_significance>
        </concept>
    <concept>
        <concept_id>10010405.10010476.10003392</concept_id>
        <concept_desc>Applied computing~Digital libraries and archives</concept_desc>
        <concept_significance>300</concept_significance>
        </concept>
 </ccs2012>
\end{CCSXML}

\ccsdesc[300]{Social and professional topics~Socio-technical systems}
\ccsdesc[300]{Social and professional topics~Technology and censorship}
\ccsdesc[300]{Social and professional topics~Governmental regulations}
\ccsdesc[500]{Applied computing~Law, social and behavioral sciences}
\ccsdesc[300]{Applied computing~Computing in government}
\ccsdesc[300]{Applied computing~Digital libraries and archives}
\ccsdesc[300]{Human-centered computing~Collaborative and social computing~Collaborative and social computing theory, concepts and paradigms}

%%
%% Keywords. The author(s) should pick words that accurately describe
%% the work being presented. Separate the keywords with commas.
\keywords{data governance; information commons; algorithmic accountability; platform governance; regulation; socio-technical systems; infrastructures; archives; municipal data; gig work; climate change}

\received{20 February 2007}
\received[revised]{12 March 2009}
\received[accepted]{5 June 2009}

%%
%% This command processes the author and affiliation and title
%% information and builds the first part of the formatted document.
\maketitle

% \section*{The Post-A(P)I Age and Three Dynamics}

For more than a century, democratic societies have depended on data to render public life visible and governable. Official statistics, regulatory filings, and administrative records have underpinned projects as diverse as labor protection, civil rights enforcement, environmental regulation, and economic planning. In each domain, new forms of measurement were accompanied by struggles over access, interpretation, and control: who could see the numbers, who could challenge them, and who was accountable for their consequences. The rise of large-scale digital platforms in the early twenty-first century appeared, for a time, to extend this evidentiary infrastructure into the everyday fabric of social interaction. 

For roughly a decade between 2008 and 2018, companies such as Twitter, Facebook, Reddit, and Google invested in APIs, data dumps, and other interfaces that created a ``golden age'' for studying social behavior at unprecedented scale and fidelity. These infrastructures functioned like scientific instruments: observatories producing high-resolution behavioral data about social structure, linguistic dynamics, cultural practices, and algorithmic interventions. Researchers in data science, computational social science, human–computer interaction, and related fields built methods, careers, and new subdisciplines on these platforms' data~\cite{Bruns_APIcalypsesocialmedia_2019,lazer_ComputationalSocialScience_2020,proferes_StudyingRedditSystematic_2021,Hampton_StudyingDigitalDirections_2017,GolderDigitalFootprintsOpportunities2014,Lazer_DataexMachina_2017}.

Platforms initially welcomed this work because it conferred legitimacy: analyses that evaluated new features, documented rapid growth, and demonstrated ``real-world impact'' made platforms more attractive to users, employees, investors, and regulators. But as these firms matured into publicly traded corporations, their incentives shifted. By the mid-2010s, researchers were using the same data infrastructures to uncover bots, harassment, polarization, disinformation, and other phenomena implicating platforms' governance failures and systemic biases.\cite{marwick_and_lewis,Walker_disinformationlandscapelockdown_2019} In response, companies adopted technical and legal strategies that made it increasingly difficult to study their systems. Access policies tightened, terms of service hardened, and claims about protecting user privacy became convenient justifications for shutting down data access to independent investigators.\cite{confessore_cambridge_2018,bobrowsky_FacebookDisablesAccess_2021,stokel-walker_under_2024} Deen Freelon aptly described the resulting condition as a ``post-API age,'' in which researchers' access to digital trace data can be restricted at any time, for any reason, and without recourse~\cite{Freelon_ComputationalResearchPostAPI_2018}. Researchers' loss of data access has only intensified since Freelon's 2018 warning. 

The emergence of large language models in the early 2020s accelerated this transformation. User-generated content accumulated over decades suddenly became highly valuable as training data. Platforms moved swiftly to prevent scraping, restrict APIs, and paywall archives, treating their data as proprietary assets rather than quasi-public resources.\cite{isaac_RedditWantsGet_2023,roose_data_2024} At the same time, user migration to alternative platforms and federated networks has made even basic measurement of online social behavior more difficult.\cite{nicholson_truthiness_2024} Few organizations outside the largest technology companies have the technical and financial capacity to sustain open data infrastructures at scale. In retrospect, the ``API age'' was not only unusually open but also riddled with inefficiencies, inequalities, and ethical abuses that platforms have seized upon to justify restricting access~\cite{boydCriticalQuestionsBig2012,Puschmann_endwildwest_2019}.

It is as if the CERN supercollider were closed to particle physicists or the James Webb Space Telescope were put behind a paywall for  astronomers. There is no straightforward path back to the peculiar data ecosystem of the late 2000s and early 2010s. Unlike particle physicists and astronomers, computational social scientists and data scientist working in the public interest generally do not design or control the instruments that generate the data on which their work depends. The challenge, however, extends beyond academic research and into the agendas of other public interest watchdogs. As data-driven systems mediate communication, work, health, mobility, and political expression, restrictions on access and scrutiny jeopardize democratic oversight more broadly.

This article argues that we should understand this situation not as a temporary policy problem but as the outcome of three deeper, mutually reinforcing forces: enclosure, exemption, and erosion. Enclosure refers to the privatization or withdrawal of data that once had public or research value, whether through API shutdowns, paywalls, or the quiet retirement of government datasets. Exemption denotes the ability of powerful platforms and AI developers to operate outside meaningful regulatory scrutiny, invoking privacy, intellectual property, or competitive sensitivity to avoid disclosure and accountability. Erosion captures the decay of institutions and infrastructures like APIs, archives, documentation practices that once enabled long-term observation and reproducible analysis of digital life. Together, these forces are transforming the informational foundations of public knowledge and weakening society's capacity to govern data-intensive systems.

This article proposes ``public interest data science'' as an interdisciplinary project grounded in explicit normative values that takes these structural pressures seriously and seeks to build countervailing commitments of openness, oversight, and ownership. It follows arguments in other domains that locate professions within longer lineages of public interest traditions like journalism, law, accounting, engineering, planning that have developed mechanisms to ensure that specialized expertise serves the common good rather than only private clients~\cite{washingtonDefiningPublicInterest2024}. Drawing on these lineages and on adjacent conversations in civic technology, data-for-good, critical data studies, and data justice, public interest data science aims to reconstruct the conditions under which data-driven systems can credibly claim to operate in the public interest.

The remainder of the article proceeds in four steps. First, it offers a conceptual diagnosis of enclosure, exemption, and erosion as structural forces undermining public oversight of data-driven systems, situating them in historical perspective alongside earlier struggles over access to statistics, financial disclosures, and environmental data. Second, it defines public interest data science and elaborates openness, oversight, and ownership as countervailing pillars, distinguishing this project from adjacent domains such as data for good, civic tech, and data justice. Third, drawing on a systematic, retrieval-augmented ``grid search'' across the literature on public interest traditions (journalism, law, accounting, engineering, planning) and public interest technology conversations (civic tech, tech for good, critical data studies), it identifies four meta-cluster: ``Evidence and Rights,'' ``Metrics and Audits,'' ``Infrastructures and Safety,'' ``Interfaces and Narratives.''

Finally, the article develops three speculative case studies to illustrate how public interest data science might be enacted in practice: mining municipal archives, infrastructuring worker observatories, and Anthropocene data trusts. These cases show how openness, oversight, and stewardship can be concretized in familiar but contested settings: local government transparency, platform-mediated work, and climate-relevant memory. The concluding discussion synthesizes anticipated critiques, draws out implications for regulation, professional norms, and grassroots practice, and sketches a forward-looking agenda for building data futures oriented toward the public, workers, and future generations.

% What then is to be done about the disciplines that depend on these data, the platforms that have become vital infrastructure for our social lives, and the myriad forces that make data paradoxically ever greater in number and also scarcer to find? Scientists should create new models for data access, new methods for studying online behavior, and new institutions capable of sustaining the kinds of research that democratic societies increasingly require. There are ample precedents to be found in the commitments other disciplines and professions have made to the public interest in the aftermath of scandals and collapses. A \textit{public interest data science} is one alternative for designing new institutional and infrastructural assemblages capable of ensuring democratic oversight of essential knowledge infrastructure. Undertaking such an effort first requires an accurate diagnosis of the forces undermining the public interest.

\section*{Three pressures}

\paragraph*{Enclosure}
The first defining force is enclosure, a term with deep historical roots. Beginning in medieval and early modern England, the enclosure movement transformed collectively held lands into private property by fencing off fields, revoking customary rights, and displacing communities who relied on shared resources. Enclosure was not merely an economic shift; it was a political act that redistributed power, restricted access, and redefined who could benefit from what had once been a commons. A similar dynamic now characterizes the digital ecosystem~\cite{boyle2017second,hess2003ideas}. Information enclosure refers to the privatization and withdrawal of access to data that once supported independent judgments about how digital systems operate. In rapid succession, Facebook restricted access to its public APIs after the Cambridge Analytica scandal, Twitter revoked research access following its 2022 acquisition, and Reddit shut down the Pushshift infrastructure while imposing steep fees for API access~\cite{isaac_RedditWantsGet_2023}. These actions mirror historical enclosure: platforms have fenced off user-generated content that previously functioned as a public or quasi-public resource for researchers, journalists, policymakers, and civil society~\cite{wu_PlatformEnclosureHuman_2021}. Governments have likewise reduced or eliminated public datasets on policing, environmental quality, public health, and education. 

\paragraph*{Exemption}
The second force is exemption, which describes how companies, their platforms, and their algorithms operate outside meaningful regulatory scrutiny. As their influence over communication, commerce, and employment has grown, these firms have invoked legal and technical shields that insulate them from oversight. Companies routinely claim privacy, intellectual property, national security, or competitive sensitivity to justify withholding data and treating opacity as a competitive advantage~\cite{ingram_silicon_2019}. They selectively invoke decades-old laws such as the Computer Fraud and Abuse Act to prevent independent data collection and Section 230 to disclaim responsibility for harms facilitated by their systems. Platforms increasingly engage in ``privacy-washing'' by claiming to protect users while in practice restricting independent research and accountability. Even in the face of well-documented harms like bias, harassment, misinformation, discrimination, and environmental externalities, companies face minimal obligation to disclose relevant data or to mitigate harms. Exemption ensures that socio-technical systems with profound social, economic, and political impact operate with limited accountability and evade responsibility. It thus creates a structural mismatch: the systems with the greatest public consequences are subject to the weakest oversight.

\paragraph*{Erosion}
Erosion, the third defining force, refers to the ongoing breakdown of the institutions and infrastructures that once enabled sustained, independent observation of platform-mediated behavior. As researchers lose access to centralized APIs and archival datasets, studying long-term trends in digital behavior becomes increasingly difficult. The rise of decentralized and federated platforms further disperses data across fragmented, incommensurable, and ephemeral infrastructures. Basic indicators that once grounded empirical research like user populations, content flows, community dynamics have become harder to measure consistently. Erosion undermines the reproducibility, comparability, and accumulation of knowledge that empirical research depends upon. It also limits the ability of journalists, policymakers, and educators to understand how digital systems shape public discourse and social life. Together, enclosure, exemption, and erosion represent a structural transformation in the foundations of public knowledge. They challenge assumptions underpinning established research disciplines and threaten the capacity of democratic societies to oversee vital systems.

%This introductory section sets the stage for examining how public interest data science draws on earlier public interest professions, distinguishes itself from adjacent fields, and contributes to the design of new regulatory and infrastructural models capable of sustaining democratic knowledge in an era of rapidly changing technological power.

\section*{Three values}
Data science is both a discipline and a profession defined broadly to include overlaps with computer and information science, statistics, economics, operations research, accounting, and geography, as well as responsibilities associated with titles such as data engineer, data scientist, business analyst, quantitative researcher, and user researcher. The focus on ``data science'' rather than other concepts like ``informatics'' or ``cybernetics'' is warranted because data science as a term is both institutionally influential and still undergoing the kind of professionalization that shaped earlier public interest fields.

At its most pragmatic level, data science is the identification of patterns and generation of explanations from complex datasets. Despite the best efforts of its boosters to emphasize neutrality and objectivity, data is inescapably political and material. As a political project, data science is embedded within a stack of practices, tools, institutions, and ideologies that shape power relations going into and coming out of data-driven systems. As a sociomaterial project, data science is the product of infrastructures, standards, and labor that collect, store, maintain, transmit, and compute information~\cite{thomerRelationalDataParadigms2020}. 

Precisely because data science has these political and material features, it is a generative site for translating the traditions and accomplishments of established public interest professions and adapting them to the challenges posed by contemporary data-driven systems. Public interest data science is a response to the transformation of the digital information environment, in which once-stable infrastructures for understanding public life are increasingly threatened by enclosure, exemption, and erosion. Against these forces, a constructive agenda grounded in openness, oversight, and ownership is needed to rebuild the conditions for democratic oversight and accountability.

\paragraph*{Openness}
Openness stands in opposition to enclosure. In a period when platforms withdraw support for APIs, governments sunset datasets, and corporations privatize user-generated content, openness asserts that the ability to observe and understand sociotechnical systems is a public necessity. This commitment involves investing in redundant archives, transparent pipelines, shared community data trusts, and durable public repositories capable of preserving access to public-relevant data even as private incentives shift. Openness also emphasizes methodological legibility: data provenance, documentation, and reproducible workflows that make it possible for more than experts to evaluate how data is collected, transformed, and interpreted. By building infrastructures that resist the fencing-off of digital commons, openness helps restore the empirical foundations of public oversight and keeps public interest lineages in journalism and law viable.

\paragraph*{Oversight}
Oversight confronts the force of exemption by insisting that data-driven systems must be governed with the same rigor, transparency, and public responsibility expected of other critical infrastructures. Platforms and AI developers increasingly claim exceptional status, evading scrutiny by appealing to the proprietary nature of their models or the supposed risks of disclosure. Public interest data science advances counter-practices: robust auditing methods, documentation standards, impact assessments, and regulatory frameworks capable of testing and verifying algorithmic claims. Drawing on institutional lessons from financial auditing, product safety regulation, and environmental monitoring, oversight treats data-driven systems as subjects of public concern rather than private fiefdoms. In doing so, it establishes the conditions under which algorithmic processes can be evaluated, contested, and governed.

\paragraph*{Ownership}
Ownership serves as the counter-force to erosion. As data disappears, infrastructures decay, and fragmented platforms undermine reproducibility, public interest data science emphasizes the importance of resilient, public stewardship of data resources. Ownership here does not refer to private possession but to collective governance, in which communities, researchers, public institutions, and civil society organizations participate in establishing norms for long-term data preservation and use. This involves building interoperable infrastructures, investing in reproducibility frameworks, creating archival strategies that safeguard public memory, and supporting institutions capable of maintaining data resources over decades. By cultivating forms of ownership that distribute responsibility and ensure continuity, public interest data science preserves the empirical record essential for democratic accountability and intergenerational justice.

\section*{Public interest lineages}
The idea of public interest data science does not emerge in a vacuum. Many professions have cultivated explicit public interest missions by developing norms, institutions, and ethical frameworks to ensure that specialized expertise serves the common good rather than private power~\cite{stapletonWhoHasInterest2022,dantecInfrastructuringFormationPublics2013,costanza-chockDesignJustice2020}. Journalism, law, accounting, planning, and engineering each grappled with harms created by opaque systems, unregulated markets, or concentrated authority. Their responses offer important lessons for public interest data science, which now faces analogous challenges in the realm of data and algorithms.

\paragraph*{Journalism}
Journalism is rooted in the democratic ideal that an informed public is necessary for self-governance~\cite{ettemaJournalismReasongivingDeliberative2007,deuzeWhatJournalismProfessional2005}. Over time, journalism developed a watchdog role: scrutinizing government, exposing corruption, and making complex information legible to citizens. Legal protections for press freedom and norms such as verification, independence, and accountability emerged to institutionalize these commitments. Investigative reporting made extensive use of public records, freedom-of-information laws, interviews, and data analysis to uncover hidden harms~\cite{glasserInvestigativeJournalismMoral1989,jacquetteJournalismEthicsTruthTelling2009,protessJournalismOutrageInvestigative1992}. Journalism's public interest tradition demonstrates that access to information is a democratic requirement~\cite{zengGhostsNewspapersPublic2022}. For data science, the parallel is clear: without durable rights to observe, analyze, and publish information about digital systems, the public cannot meaningfully evaluate their impacts.

\paragraph*{Law}
Law developed its own public interest orientation through legal aid movements, civil rights litigation, pro bono service, and cause lawyering~\cite{council1976balancing,rhode2004probono,sarat1998cause,cummings2009public}. Public interest lawyering emerged in response to structural inequality, recognizing that formal rights are meaningless without institutions that help people exercise them~\cite{southworth2005what,cummings2006after,trubek2011public}. Lawyers working in the public interest developed strategic litigation, community partnerships, and impact-driven advocacy to challenge discriminatory policies, regulate dangerous industries, and expand civil liberties~\cite{cummings2009public,nielsen2006organization,albiston2017public}. They also designed institutions such as public defender systems, legal clinics, and nonprofit advocacy organizations that could sustain this work outside market pressures~\cite{rhode2004probono,albiston2017public}. Public interest data science must similarly confront problems that are not neutral or evenly distributed: just as public interest lawyers challenge powerful actors, public interest data scientists must illuminate opaque algorithmic systems, quantify disparate impacts, and support communities seeking redress~\cite{barocas2016big,selbst2019fairness}.

\paragraph*{Accounting}
Accounting provides one of the clearest historical analogies~\cite{carnegieCriticalInterpretiveHistories1996,leeProfessionalizationAccountancyHistory1995}. In the early twentieth century, recurrent financial crises revealed that private auditing practices routinely failed to protect investors or the public. In response, the profession articulated a explicit "public interest mandate": auditors were to serve not the firms that hired them, but the broader public whose well-being depended on reliable financial information~\cite{quayleWhistleblowingAccountingPublic2021,bakerWhatMeaningPublic2005,dellaportasReflectionsPublicInterest2008}. Regulatory institutions such as the Securities and Exchange Commission (SEC) established standards for disclosure, independence, and verification. Algorithmic systems resemble corporate balance sheets that can hide risks, mask errors, or facilitate fraud unless subject to independent audit. Public interest data science can draw directly from accounting's lessons to design methods for algorithmic auditing, documentation, stress-testing, disclosure, and transparency that serve public interests rather than the entities deploying the systems.

\paragraph*{Planning}
In twentieth-century cities, urban planning practices were often technocratic and exclusionary, producing highways, zoning policies, and redevelopment projects that disproportionately harmed marginalized communities~\cite{jacobs1961death,doryClashUrbanPhilosophies2018,guttenbergPlanningIdeology2009}. The advocacy planning movement emerged in response, insisting that planners had obligations to represent underserved populations, engage impacted communities, and make planning processes more participatory~\cite{dadashpoorDefiningPublicInterest2021}. Planning has also shifted toward more holistic and ecological perspectives that recognize problems such as traffic, affordability, and climate risk as deeply entangled~\cite{hirtPremodernModernPostmodern2009,roheLocalGlobalOne2009}. Planning's public interest tradition highlights that expertise must be paired with democratic responsiveness and an awareness of power~\cite{lanePublicParticipationPlanning2005}. For public interest data science, the parallel is the need for participatory design, community stewardship of data, and recognition that technical decisions about what data to collect, what models to build, and how to interpret results have distributive consequences.

\paragraph*{Engineering}
Engineering likewise evolved a public interest ethos, especially in response to infrastructure failures, industrial accidents, and technological risks in the nineteenth and twentieth centuries~\cite{petroski1992engineer,perrow1984normal,vaughan1996challenger,dekkerDriftFailureHunting2016}. Engineering societies developed codes of ethics requiring practitioners to hold paramount the safety, health, and welfare of the public alongside adopting practices of risk assessment~\cite{nathwaniThreePrinciplesManaging1995,boudiaIntroductionRiskRisk2007}. Regulatory bodies emerged to ensure that bridges, buildings, railroads, cars, pharmaceuticals, and medical devices met rigorous safety standards through testing, certification, and oversight of occupational health and safety~\cite{abramsShortHistoryOccupational2001}. Managing complex systems required cultural shifts to embrace practices of safety and high reliability rather than techno-solutionism~\cite{weickOrganizingHighReliability1999,haleWorkingRuleWorkingI2013,haleWorkingRuleWorkingII2013}. Public interest data science should emulate engineering's commitments to safety, standards, and oversight: data and algorithmic systems increasingly function as critical infrastructure and should be governed accordingly.

\section*{Public interest conversations}
Public interest data science shares intellectual ground with several adjacent fields but is distinct in its commitments, methods, and institutional ambitions. These neighboring traditions help frame important ethical, technical, and societal concerns around data systems, but they often approach these concerns with goals other than building durable infrastructures and regulatory regimes for public accountability.

\paragraph*{Data for social good}
Data and AI for social good initiatives emerged largely within philanthropic, corporate, and nonprofit contexts, applying data science to social challenges such as poverty, education, public health, or disaster response~\cite{williamsDataActionUsing2022}. These efforts have brought timely analytical capabilities to organizations with limited technical resources and emphasized socially beneficial applications of data~\cite{cowlsAISocialGood2021}. However, data for good is often project-based, tool-centric, and oriented toward helping institutions manage problems rather than questioning the broader systems that produce harm~\cite{floridiHowDesignAI2020}. It rarely interrogates power, governance, or structural inequities and often assumes that private-sector partners will remain benevolent stewards of data access---assumptions the post-API age has called into question~\cite{Freelon_ComputationalResearchPostAPI_2018,taylorPublicActorsPublic2021,schwoererWhoseOpenData2022,trombleWhereHaveAll2021}. Public interest data science differs in insisting that long-term accountability requires more than well-intentioned applications; it requires durable infrastructures, transparent practices, and regulatory constraints that counteract enclosure, exemption, and erosion across the data ecosystem.

\paragraph*{Civic technology}
Civic technology shares a complementary but distinct lineage. It focuses on building digital tools that improve public service delivery, enable citizen engagement, or streamline administrative processes~\cite{boehnerDataDesignCivics2016,mcguinnessPowerPublicPromise2021}. Civic tech has played a vital role in modernizing government IT systems, facilitating new forms of public participation, and creating new institutions and professional identities~\cite{aragonCivicTechnologiesResearch2020,knodelShapingProfessionBuilding2025,riderBuildingIdealWorkplaces2022}. Yet civic technology often prioritizes efficiency and practical problem-solving without addressing deeper questions of data rights, algorithmic accountability, or the political economy of digital infrastructures~\cite{aulaSteppingBackData2023,chambersBigTechAdvocacy2025,stapletonWhoHasInterest2022}. It also tends to operate within existing institutional arrangements rather than challenging them~\cite{chowdharyCareLiberationCreating2020,milanAlternativeEpistemologiesData2016a}. Public interest data science, by contrast, sees civic capacity not simply as a matter of better tools but as a matter of reforming or designing alternative institutions for archiving, oversight, and stewardship to ensure accessibility and accountability.

\paragraph*{Critical data studies}
Critical data studies offers an essential set of theoretical and interpretive tools, examining how data practices shape categories of identity, produce forms of surveillance, and reinforce existing power relations~\cite{boydCriticalQuestionsBig2012,dignazioDataFeminism2020}. Scholars in this tradition have illuminated the political, historical, and ethical stakes of datafication, calling attention to issues such as classification violence, algorithmic inequity, and the social construction of data~\cite{burgessEverydayDataCultures2022,greenDataSciencePolitical2021,taylorWhatDataJustice2017}. Public interest data science stands on these shoulders, especially in understanding how enclosure, exemption, and erosion threaten democratic knowledge. Yet critical data studies often stops at critique and is less explicit about designing alternative infrastructural or regulatory futures~\cite{neffCritiqueContributePracticeBased2017}. Public interest data science extends the critical tradition by translating analysis into institutional commitments and design principles grounded in precedents from other public interest professions.

\paragraph*{Data justice}
Data justice, originating in both grassroots activism and academic inquiry, centers equity, rights, and the distribution of harms and benefits in data systems. It foregrounds the experiences of communities disproportionately affected by discriminatory algorithms, predatory data collection, and harmful data governance practices. Public interest data science shares these normative commitments and seeks to make them operational within broader governance structures. But while data justice focuses on justice-oriented outcomes, it does not always engage the technical, archival, or administrative infrastructures required to ensure ongoing accountability. Public interest data science therefore positions itself as the bridge between justice-oriented values and the institutional systems like regulatory regimes, public data repositories, transparency requirements, auditing practices, and participatory governance necessary to realize those values at scale.

\section*{Anticipating Critiques}

Any agenda that tries to reshape the political economy of data will encounter sustained, and often principled, resistance. Public interest data science challenges familiar arrangements of data access, corporate secrecy, regulatory design, and epistemic authority. It asks platforms to open systems they currently control, regulators to govern domains they only dimly understand, and technical experts to treat public accountability as a core part of their role. It is therefore unsurprising that industry leaders, policymakers, academics, and even sympathetic critics raise objections. Taking those critiques seriously is not a cosmetic exercise. It is an opportunity to clarify what public interest data science is—and is not; to distinguish it from adjacent efforts; and to refine its methods and ambitions. Below, I group the most common objections into six families—realist, critical, libertarian, positivist, pessimist, and pragmatist—and show how engaging each one helps sharpen the field's agenda.

\subsection*{Realist}
\begin{quote}\centering
\textit{``The open data era is over and never coming back.''}
\end{quote}

\noindent The realist objection begins from an uncomfortable truth: the conditions that made the “API era” possible have changed. Platforms that once courted developers and researchers with generous APIs now treat data as a scarce strategic asset. Governments have quietly withdrawn or underfunded key public datasets. From this vantage point, calls for renewed openness can look nostalgic: enclosure appears as the natural maturation of technological and economic systems, not a reversible choice.

Public interest data science accepts the descriptive force of this story but rejects its fatalism. Enclosure is real, but not destiny. Historically, sectors far more opaque and powerful than today’s platforms have been opened to scrutiny. Early financial markets operated with minimal disclosure until fraud and crises triggered securities law; polluting industries hid their impacts until monitoring and reporting requirements made emissions visible; pharmaceutical firms resisted trial transparency until regulators and advocates forced reporting standards. None of these shifts resulted from voluntary altruism. They emerged when publics, journalists, lawyers, and regulators constructed new regimes of transparency and oversight.

Realism is valuable when it exposes the limits of corporate goodwill; it becomes a problem when it slides into resignation. Public interest data science reframes enclosure as a political condition that can be contested through statutory access rights, mandated research portals, publicly funded observatories, and hybrid governance models—turning realist critique into an argument for institutional reform rather than an excuse for inaction.

\subsection*{Critical}
\begin{quote}\centering
\textit{``How is this different from tech for good, civic tech, or data justice?''}
\end{quote}

\noindent From data justice, critical data studies, and long-standing public interest technology communities comes a different concern: proliferation fatigue. For those who have spent years critiquing data harms and building alternatives, “public interest data science” can sound like yet another label layered atop data for good, civic tech, responsible AI, and more.

This critique matters because it forces clarity about distinctiveness. The claim is not that public interest data science is wholly new, but that it integrates and redirects existing strands toward institutional change. Data-for-good initiatives often center projects—hackathons, pilots, philanthropic collaborations—that rarely alter underlying power structures. Civic tech emphasizes tools and services, frequently aimed at improving government performance rather than reshaping data governance itself. Critical data studies and data justice provide essential analyses of race, class, gender, and colonialism and powerful normative visions, but are not always oriented toward building durable sociotechnical systems.

Public interest data science self-consciously sits at this intersection. Its wager is that we need an agenda that holds together structural critique, technical practice, institutional design, and professional identity. The critical objection thus becomes a constructive challenge: do not rebrand; coordinate. Unless public interest data science invests in concrete pathways—public data agencies, audit bodies, archives, and professional norms—it will rightly be dismissed as just another slogan.

\subsection*{Libertarian}
\begin{quote}\centering
\textit{``Oversight will kill innovation.''}
\end{quote}

\noindent The libertarian critique reprises a familiar argument: tighter oversight and disclosure will slow innovation, deter investment, and erode competitiveness. In the data and AI context, this appears as warnings that audit mandates, transparency rules, or data-access obligations will chill experimentation and burden startups more than incumbents. There are real design challenges here—poorly crafted regimes can entrench large firms and overwhelm smaller actors. But the broader claim that accountability and innovation are inherently at odds does not withstand historical scrutiny. Industries once hostile to regulation now rely on it to sustain public trust. Commercial aviation did not flourish in spite of safety standards; it flourished because passengers trusted that planes were safe. Financial reporting requirements impose costs but help prevent crises that would destroy markets. Environmental regulations have often spurred cleaner technologies and new sectors rather than freezing progress.

Public interest data science draws on these analogies to argue that well-designed oversight can be pro-innovation when innovation is understood as long-term, socially beneficial progress rather than short-term arbitrage. Documentation requirements, incident reporting, and independent audits can level the playing field, reward robustness and fairness, and reduce the risk of catastrophic failures that invite blunt bans. The libertarian critique nevertheless usefully underscores that regulatory design matters: accountability mechanisms must be iterative, evidence-informed, and co-developed with diverse stakeholders to avoid unintended harms.

\subsection*{Positivist}
\begin{quote}\centering
\textit{``Data science should not be political.''}
\end{quote}

\noindent From within technical disciplines comes a different concern: that explicitly orienting data science toward “public interest” values risks politicizing a field that should remain neutral, objective, and method-driven. In this positivist critique, algorithms are tools; ethics should reside in downstream use, not in methods themselves. If every modeling decision is cast as political, the worry goes, we invite relativism and paralyze practice.

Public interest data science treats this as a chance to clarify what “political” means. The claim is not that every modeling choice is partisan, but that no pipeline is value-free. Selecting target variables, features, loss functions, and thresholds encodes assumptions about what counts as success, which errors are acceptable, and whose experiences matter. A recommender that optimizes for engagement privileges some attention and content; a risk model tuned to minimize one type of error will shift burdens onto particular groups.

Making such choices explicit does not undermine rigor; it strengthens it. It invites documentation, multidisciplinary scrutiny, and contestation by affected communities, aligning with emerging practices like datasheets, model cards, and participatory design. The positivist critique is useful to the extent it demands concrete, operational commitments rather than vague moralizing. Neutrality, in this view, is not the absence of values but honesty about which values are at work and how they are implemented in code, data, and evaluation.

\subsection*{Pessimist}
\begin{quote}\centering
\textit{``These systems are too complex to oversee.''}
\end{quote}

\noindent The pessimist critique centers on the opacity of modern data systems. Deep learning models with billions of parameters, constantly updated recommender engines, and sprawling data pipelines are difficult to understand even for their creators. If the systems themselves are complex and rapidly evolving, skeptics ask, how can external auditors, regulators, or the public meaningfully oversee them? Calls for accountability can sound naïve against this background, or collapse into vague demands for “full transparency” that are technically and organizationally infeasible.

Public interest data science reframes this as a design problem, not a conversation stopper. We already govern complex systems—nuclear power, air traffic, high finance—through layered oversight, specialized roles, and structured disclosure. No single person grasps every detail of a jet engine or derivatives book; instead, we build standards, redundancy, and monitoring regimes that make catastrophic failures rarer and more manageable.

In the algorithmic realm, this implies designing for “auditable complexity.” Some oversight targets inputs and processes: documentation of training data, change logs, records of design choices. Other oversight targets outputs and impacts: stress tests, performance by subgroup, incident investigation. Explainability tools, proxy metrics, and rigorous behavioral evaluations all play a role. The pessimist critique is valuable insofar as it pushes public interest data science beyond slogans toward graduated access, role-based visibility, and independent intermediaries—akin to financial regulators and auditors—who can see more than the public while remaining accountable to it. Complexity raises the bar for oversight; it does not eliminate the need.

\subsection*{Pragmatic}
\begin{quote}\centering
\textit{``Resources, privacy, and public interest fatigue''}
\end{quote}

\noindent The pragmatic critique highlights real constraints: building public data infrastructures, observatories, and oversight bodies is expensive; opening data can create privacy risks; and sustaining public engagement with technical governance is difficult when attention is scarce. Public interest data science takes these limits seriously. Accountability is not free—it requires funding, expertise, and time—but so do courts, schools, and roads. If data-driven systems now shape credit, employment, health, and information access, then investing in their governance is part of basic civic infrastructure. The costs of unchecked systems—discrimination, disinformation, systemic error—are often far higher.

Privacy is likewise a genuine constraint. Calls for openness can clash with the need to protect people from surveillance or re-identification. The answer is not to abandon transparency but to design it carefully: aggregation, differential privacy, tiered access, robust consent, and strong legal safeguards. Openness and privacy are not opposites; they are jointly necessary for legitimate data governance.

Finally, not everyone will engage with audit reports or data schemas. This is why public interest data science emphasizes interfaces, narratives, and intermediaries like journalists, advocates, librarians, professional associations that translate technical oversight into accessible explanations and action. The pragmatic critique usefully insists that proposals be implementable, trade-offs explicit, and benefits visible, helping prioritize where openness, oversight, and stewardship will matter most.

\section*{RAG-augmented thematic grid search and meta-analysis}
To situate public interest data science within existing professional and technological lineages, I employed a structured, explicitly normative ``grid search'' and meta-synthesis over a two-dimensional conceptual space. One axis comprised public interest traditions (PIT)—journalism, law, accounting, engineering, and planning—each with established public mandates and institutional histories. The other axis comprised public interest technology conversations (PIC): civic tech, tech for good, critical data studies, and data justice. The resulting 5×4 matrix produced twenty PIT×PIC pairings, treated as individual ``cells'' for systematic theoretical exploration.

The first step was to operationalize each PIT and PIC in terms of their characteristic values, practices, infrastructures, and vocabularies. For example, public interest journalism was characterized by watchdog functions, open records, and narrative accountability; law by rights-claiming, litigation, and institutional remedies; accounting by auditing, materiality, and financial disclosure; engineering by safety, standards, and codes of ethics; and planning by participatory design, spatial justice, and long-term public goods. On the PIC side, civic tech emphasized open interfaces, co-production with governments and communities, and practical tools; tech for good foregrounded philanthropic projects, impact narratives, and solutionist logics; critical data studies stressed power, classification, and structural critique. This scoping generated a shared vocabulary for comparing and integrating traditions.

For each of the twenty PIT×PIC cells, I then produced an approximately 250-word narrative synthesis addressing four questions: (1) What shared values or practices make this pairing promising for public interest data science? (2) Where do their norms or temporalities conflict (e.g., slow legal processes vs. fast-moving platforms)? (3) What are the implications for governing data and algorithmic systems as commons? (4) What concrete opportunities for collaboration or institutional design emerge from this pairing? These narratives were deliberately written as generative hypotheses, sketching both complementarities and frictions rather than simply affirming alignment.

Each narrative was then scored for persuasiveness on a 1–10 scale against an explicit normative prompt: How compelling is this pairing as a foundation for public interest data science oriented toward democratic control, commons governance, and accountability of data-driven systems? This scoring was not treated as a quantitative measure of ``truth,'' but as a disciplined way to prioritize which hybrid lineages appeared most fertile for further development. The explicit normative basis—anchored in public interest, commons stewardship, and democratic oversight—ensured that the evaluation criteria were transparent rather than implicit.

In a second pass, I generated parallel matrices of strengths, weaknesses, tensions, and opportunities for each PIT×PIC cell. Strengths captured the most promising shared capacities (e.g., journalism×critical data studies combining investigative practice with structural critique). Weaknesses identified internal limitations or blind spots (e.g., tech for good's reliance on philanthropic funding and pilotism). Tensions focused on structural frictions (e.g., engineering's risk-averse standardization versus civic tech's experimentalism). Opportunities articulated plausible joint projects, institutions, or infrastructures (e.g., law×civic tech supporting FOIA-driven algorithmic disclosure tools). Each of these four dimensions was summarized as a one-line statement per cell, yielding a compact comparative view of pros, cons, tensions, and consensus opportunities.

To deepen the analysis, I then generated ``opposing statements'' for each synthesis, articulating skeptical or countervailing readings of the pairing (for example, how journalism×tech for good could devolve into PR-driven impact storytelling rather than adversarial investigation). These opposing statements were also scored for persuasiveness against the same normative prompt. Comparing scores between base and opposing narratives helped identify where the case for a given PIT×PIC pairing was robust versus fragile, and where tensions might fundamentally limit its usefulness as a building block for public interest data science.

With these textual artifacts in hand—twenty synthesis narratives, four one-line matrices (strengths, weaknesses, tensions, opportunities), and their oppositional counterparts—I treated each cell as a conceptual ``document'' embedded in a latent semantic space. Rather than performing a purely computational embedding, I conducted an interpretive clustering exercise grounded in recurring themes, metaphors, and institutional functions. I looked at which cells repeatedly invoked evidence, rights, and watchdog roles; which centered metrics, audits, and reporting; which emphasized infrastructures, space, and safety; and which revolved around interfaces, narrative, and public understanding. This step effectively sampled the conceptual space spanned by the PIT×PIC grid and grouped cells by thematic proximity rather than by their row or column positions.

Through iterative grouping and re-grouping—while enforcing the design constraint that clusters should mix different PIC types rather than simply reproduce existing categories—I converged on four meta-clusters:

\subsection*{Meta-cluster 1: Evidence, Watchdogs, and Rights}
%(e.g., journalism×civic tech, law×critical data studies)
This cluster frames data as a tool of evidence to uphold rights, inform policy, and advance justice. It draws heavily from the legacy of journalism and law. Just as investigative journalists leverage public records and freedom-of-information laws to expose hidden harms, public interest data scientists work to ensure the right to access information about sociotechnical systems. Transparency is not an abstract ideal here but a precondition for rights – without durable rights to observe, analyze, and publish information about algorithms and platforms, the public cannot evaluate or contest their impacts. This frame includes efforts to establish data access as a civil right (for example, treating certain platform data as subject to public records laws or scholarly access regimes) and to use data analysis as evidence in service of civil rights and human rights. Think of cases where data science helps prove discrimination (e.g. showing racial bias in lending algorithms or policing strategies) and thereby supports legal action or policy change. ``Evidence and Rights'' also means empowering communities with data that substantiate their claims – for instance, helping a neighborhood gather environmental data to enforce pollution regulations, or enabling voters to audit public spending. This cluster resonates with the public interest law tradition that formal rights mean little without institutions to enforce them. Thus, it involves building those institutions (data clinics, legal data trusts, transparency portals) that make data evidence available for advocacy and accountability. In summary, Evidence and Rights orient data science toward supporting an informed citizenry and an equitable legal order, ensuring that the empirical basis for justice is robust and accessible.


\subsection*{Meta-cluster 2: Metrics, Audits, and Institutional Accountability}
%grouping cells oriented toward quantification as a governance device (e.g., accounting×civic tech, engineering×tech for good).
This cluster is about developing the measurement and verification tools needed to hold algorithms and data-driven processes accountable. It is inspired by the field of accounting and auditing, which over a century ago established that professionals must serve the public interest by verifying corporate disclosures. Today's analog is the independent audit of algorithms, datasets, and AI models. Metrics and audits work involves creating techniques to test algorithms for bias, error, or misuse; designing fairness and accountability metrics; and instituting regular auditing processes much like financial audits. For example, researchers might devise a ``stress test'' for a ride-hailing platform's pricing algorithm to see how it behaves under different scenarios, or an NGO might audit a government predictive policing system for racial bias. Public interest data science explicitly looks to accounting's lessons here: just as corporate balance sheets can hide risks or fraud unless independently audited, algorithmic systems can mask inequities or errors unless we subject them to scrutiny. This cluster also extends to developing standards and documentation (model cards, datasheets, audit frameworks) that make it feasible for external parties to evaluate complex AI systems. Importantly, Metrics and Audits are not just technical exercises; they require institutional mandates. For instance, regulators might require annual algorithmic impact assessments from tech firms, analogous to annual financial audits. An emerging example is the concept of an algorithmic auditor or model audit trail, echoing the rise of the public company auditor in the 20th century. Ultimately, this cluster gives substance to oversight: it provides the methodologies and professional practices by which we can verify that data-driven systems are doing what they claim, not causing undue harm, and conforming to legal and ethical standards. It's a direct answer to the question, ``How do we actually hold algorithms accountable?'' – by measuring them, testing them, and institutionalizing those verification processes as a matter of public interest.

\subsection*{Meta-cluster 3: Infrastructures, Space, and Socio-Technical Safety}
% encompassing pairings that treat data and models as part of safety-critical, spatial, and environmental infrastructure (e.g., engineering×critical data studies, planning×civic tech).
This cluster emphasizes the engineering and governance of technical infrastructure to ensure safety, reliability, and long-term public benefit. It parallels the ethos of civil engineering and urban planning, which evolved public interest commitments after historical failures and injustices in infrastructure development. In data science, Infrastructures and Safety means treating data and AI systems as critical infrastructure that must be designed and regulated with public welfare in mind. For example, consider how power grids or bridges are subject to safety codes, stress tests, and oversight bodies. Similarly, major algorithmic systems (like autonomous vehicle software or large-scale public data platforms) should meet rigorous standards for safety, robustness, and risk mitigation. This cluster encompasses work on building resilient data infrastructure – things like open-source tools, public data repositories, privacy-protective data sharing platforms, and reliable archival systems. It also connects to the idea of ``data commons'' and community-owned infrastructure: projects that create shared data resources governed by public or cooperative models rather than by profit. Safety, in this context, is both technical (preventing accidents, failures, data breaches) and societal (preventing systemic risks like reinforcing oppression or destabilizing institutions). For instance, ensuring the safety of an algorithm might involve external red-team testing for worst-case scenarios (similar to how new drugs are tested for side effects), or stress-testing a civic open data portal to guarantee it can't be easily taken down or corrupted. It also means having contingency plans and redundancies – backing up important datasets in multiple locations (mirrored archives) in case one fails or is sabotaged. A public interest data scientist in this cluster might work on something like an early-warning system for detecting when a critical public dataset (say, health statistics) is at risk of loss or manipulation, much like environmental engineers monitor water quality to prevent crises. The overarching philosophy is that data systems serve as public infrastructure in modern society, so they should be managed with the same care and accountability as our roads, utilities, and buildings. This entails embracing an ethic of ``paramount the safety, health, and welfare of the public'' – echoing engineering codes – in the design and deployment of algorithms and databases that people now depend on.

\subsection*{Meta-cluster 4: Interfaces, Narratives, and Data Power}
%bringing together pairings concerned with how people encounter, interpret, and contest data systems (e.g., journalism×tech for good, planning×critical data studies).
The final cluster focuses on the human dimensions of data – how data is translated into information that communities can understand and act upon. This area draws from journalism, human-computer interaction, and design. It recognizes that even the best data or analysis has limited impact if it isn't made legible and relevant to the public. Public interest data science must therefore excel at creating interfaces (visualizations, tools, platforms) and narratives (stories, explanations, framings) that connect data to democratic dialogue. This could mean developing intuitive dashboards that help citizens explore budget data, designing participatory mapping tools for neighborhoods to chart inequities, or practicing data storytelling that brings statistical insights to life through compelling narratives in media. The heritage here is clear: journalism's watchdog role and communication skill in making complex information accessible is a foundation. For instance, investigative journalists have long combined data analysis with human stories to spur action – from using spreadsheets of campaign finance records to trace corruption, to mapping data on lead poisoning to prompt a public health response. Public interest data scientists working on Interfaces and Narratives might collaborate with news organizations to produce interactive explainers about algorithms (demystifying how a particular AI system works and whom it impacts), or with community groups to develop narrative-driven reports on local data (turning a heap of public records into an accessible story about city services). This cluster also stresses participatory design: involving the public in creating the tools and narratives so that they resonate culturally and address real needs. For example, an interface for a community pollution monitoring project might be co-designed with residents to ensure it highlights the issues they care about and is easy to use. Meanwhile, narrative competence ensures that data doesn't remain in a silo of experts – it becomes part of broader social conversations. In summary, Interfaces and Narratives ensure that the insights and transparency gained through the other clusters actually reach and empower the public. By building bridges between data and democracy – whether through engaging visuals, accessible apps, or powerful stories – this cluster helps cultivate an informed public that can leverage data as a shared language for discussing and governing societal issues.

\section*{Three Cases}
Public interest data science is not solely a methodological shift but a civic and institutional project to rebuild and maintain the conditions under which democratic societies can understand, evaluate, and govern the powerful data-driven systems that increasingly shape public life. What distinguishes public interest data science is its insistence that democratic oversight requires not only ethical commitments but also material infrastructures, public institutions, and durable accountability mechanisms capable of resisting enclosure, countering exemption, and preventing erosion across the data ecosystem. Three case studies illustrate how these dynamics play out in practice and how public interest data science reconfigures the relationship between data, power, and the public good.

\subsection*{Case 1: Mining municipal archives}
% Disclosure and transparency rules shape much of what local governments can do
% But the products of this transparency is hidden away in PDFS
% Disclosure and transparency rules shape much of what local governments can do
% AI RAG and information retrieval
% AI Pulse and Valuing Human-Generated Training Data
% Clawing data back from enclosure through bad standards and implementing openness
% Which models, pipelines, and file formats
Local governments sit at the center of a long public interest tradition: open-records laws, public budgeting, zoning decisions, environmental monitoring, police oversight, school governance, and the routine administration of services. Each meeting, permit, inspection, and vote leaves a trail—council packets, land-use maps, police logs, budget ledgers, environmental reports, procurement contracts, housing documents, and voting records. Taken together, these materials form a vast, largely unstructured corpus of human-generated data about how power operates close to home.

On paper, this is a cornerstone of democratic transparency. Sunshine laws, public records statutes, and open-meetings requirements were designed to ensure that residents, journalists, researchers, and advocacy groups can see how decisions are made and resources allocated. In practice, however, the material reality of disclosure often frustrates that promise. Documents are posted as scanned PDFs, scattered across departmental websites, buried behind brittle search interfaces, or available only through formal requests that take weeks and sometimes require substantial fees. The result is an archive that is simultaneously public and opaque—disclosed yet unreachable. The records exist, but they might as well be locked in a basement filing cabinet.

This case study on ``mining municipal archives,'' starts from that tension. It asks what it would take to turn these inert, fragmented disclosures into a living evidence base for oversight and public learning. It illustrates how public interest data science can work at the intersection of law, journalism, accounting, planning, and engineering to reconstruct municipal transparency as a practice rather than a box-checking exercise—and how emerging AI tools both complicate and expand that project.

\subsubsection*{Pressures and values}
Municipal archives offer a concrete view of the three structural forces discussed earlier—enclosure, exemption, and erosion—and the countervailing commitments to openness, oversight, and ownership.

\paragraph{Enclosure vs. openness}
Enclosure here rarely takes the dramatic form of a platform shutting down an API. Instead, it accumulates through mundane decisions: saving documents as image-only PDFs, omitting metadata, reorganizing websites without redirects, or limiting search to titles rather than full text. A resident trying to trace changes in policing policy might find that relevant reports are spread across multiple sites, mislabeled, or embedded in scanned attachments that cannot be searched. Formally, the records are ``open.'' Functionally, they are enclosed by poor standards and outdated tools. Public interest data science reframes open government not as a matter of uploading files but as a commitment to usable openness: machine-readable formats, consistent schemas, and search and retrieval infrastructures that lower the cost of discovery. Openness in this sense is not a one-time act but an ongoing infrastructural responsibility.

\paragraph{Exemption vs. oversight}
Municipal records also reveal subtle forms of exemption. When transparency is treated as mere compliance—meeting minimum statutory requirements without regard to usability—local governments effectively exempt themselves from meaningful public scrutiny. Decisions about what gets recorded, how promptly it is posted, and how easy it is to find become quiet levers of power. Oversight, in a public interest frame, means more than the formal right to request documents. It requires structured observability: archives that make it possible to trace the life of a policy, follow money across budgets and contracts, and connect decisions to outcomes. Mining municipal archives with AI-enhanced retrieval and structuring tools can help build that observability, but only if those tools themselves are governed as public-interest infrastructure rather than proprietary products.

\paragraph{Erosion vs. ownership} 
Finally, municipal archives are vulnerable to erosion. Websites are redesigned; links rot; older documents are lost in migrations; file servers fail; software vendors change. Over decades, these small failures can erase crucial parts of the historical record, undermining longitudinal analysis of budgets, enforcement patterns, or land-use trends. Ownership, in the public interest sense, is about stewardship rather than control. It entails building resilient, redundant, and well-governed archival infrastructures: persistent identifiers, mirrored repositories, community backup projects, and clear custodial responsibilities. Treating municipal data as a shared civic asset—owned in common and maintained for future generations—shifts the question from ``Is this online today?'' to ``Will this still be findable and interpretable in 10 or 30 years?''

\subsubsection*{Meta-cluster mappings}
This case also illustrates how the theoretical PIT×PIC grid and its four meta-clusters. Mining municipal archives is not a single technical project but a convergence point for multiple public interest lineages: the watchdog ethos of journalism, the rights orientation of law, the audit discipline of accounting, the safety culture of engineering, and the participatory aims of planning—amplified by civic tech practices and critical perspectives on power.

\paragraph{Evidence and rights}
At its core, the municipal archive is an evidentiary commons. It documents how budgets are allocated, who receives permits, where complaints cluster, and how policies evolve. Public interest journalism and law have long relied on such records for investigations and litigation. Mining municipal archives with retrieval-augmented generation and document understanding tools strengthens this watchdog ecosystem: it becomes easier to assemble chains of evidence, compare neighborhoods, and surface patterns that would be invisible in unstructured PDFs. This aligns with the evidence and rights meta-cluster: data is not an abstract resource, but a substrate for rights-claiming and democratic oversight.

\paragraph{Metrics, audits, and institutional accountability}
Once archives are structured and linkable, they can feed new forms of public metrics and audits. For example, one could construct time series showing how often zoning variances were granted in different districts, or how inspection backlogs correlate with complaint volume. These are not just descriptive charts; they are potential audit tools that, like financial statements, allow residents and regulators to ask: Are commitments being met? Are resources distributed equitably? Public interest data science, drawing on accounting and engineering traditions, can help design these metrics so they illuminate institutional performance rather than obscure it.

\paragraph{Infrastructures, space, and socio-technical safety}
Municipal archives are also deeply spatial and infrastructural. Land-use records, transportation plans, flood maps, and building permits all shape the built environment and environmental risk. Integrating these records into interoperable data infrastructures—where decisions about roads, zoning, and climate adaptation can be analyzed together—speaks directly to the infrastructures and safety cluster. Here, mining archives is not only about past accountability but about improving future planning: enabling communities to see how infrastructure investments and hazards are distributed, and to contest inequities.

\paragraph{Interfaces, narratives, and data power}
Finally, the value of mined archives depends on how they are presented and narrated. Interfaces that allow residents to search for ``every time a planned development near my neighborhood was delayed,'' or to visualize budget shifts over time, translate raw records into actionable understanding. Storytelling—through local news, community workshops, or interactive explainers—turns patterns discovered in the data into public narratives that can inform debate. This is the interfaces and narratives cluster in action: designing touchpoints where people encounter data about their city in ways that are intelligible, contestable, and linked to civic action.

\subsubsection*{Anticipating critiques}
A serious public interest agenda around mining municipal archives has to take seriously at least three lines of critique: worries about feasibility and capacity, concerns about politics and privacy, and doubts about whether these infrastructures will genuinely redistribute power or simply entrench new intermediaries.

The first critique is about feasibility and uneven capacity. Many local governments already struggle to maintain basic IT systems, let alone robust, machine-readable archives. Staff are overextended; recordkeeping practices are inconsistent; legacy systems, vendor lock-in, and technical debt are the norm. From this vantage point, proposals to standardize formats, improve metadata, and layer AI-based retrieval on top of scattered PDF repositories can look like a form of techno-solutionism—aspirational for well-resourced cities, unrealistic for the vast majority of jurisdictions. Critics argue that without sustained funding, training, and institutional support, mining archives at scale risks producing islands of excellence surrounded by large deserts of neglect. A public interest approach must therefore treat capacity building, shared services, and regional or cooperative infrastructures as central, rather than assuming that every municipality can independently become a high-functioning data steward.

A second line of critique focuses on politics, representation, and privacy. Municipal archives are not neutral reflections of local life; they are selective records shaped by what governments choose to document and how. Mining these archives may amplify biases in what has historically been recorded—police reports and code violations may be rich, while tenant complaints or informal community knowledge are sparse. Turning more of this material into searchable, linkable data can make already over-policed or over-regulated communities even more legible to enforcement, even as other forms of inequity remain under-documented. At the same time, detailed records about individuals—addresses, complaints, health inspections—raise real privacy concerns when made more easily discoverable through powerful search and summarization tools. Critics warn that without careful redaction, access controls, and community governance, efforts to enhance transparency could inadvertently facilitate commercial exploitation, harassment, or new forms of surveillance. Public interest data science has to own this tension: making archives more usable for oversight must be paired with strong privacy safeguards and explicit attention to whose visibility is being increased, and to what end.

A third critique speaks to power and co-optation. Mining municipal archives requires tools and infrastructure; in practice, this often means relying on private vendors, cloud platforms, or third-party civic tech projects. Skeptics worry that ``opening up'' archives via commercial platforms could amount to a new round of enclosure: public records are cleaned, indexed, and served through proprietary systems, and meaningful access becomes contingent on subscriptions, terms of service, or opaque ranking algorithms. There is also the risk of ``transparency theater'': sleek dashboards and AI assistants that create an appearance of openness while leaving underlying budgeting practices, land-use politics, or enforcement priorities untouched. In this view, mining archives could produce more sophisticated interfaces for the same asymmetries—helping political actors manage narratives rather than enabling residents to contest decisions.

Taken together, these critiques do not argue against making municipal archives more usable; they specify what would be required for such efforts to be genuinely public-interest. They highlight that technical ambition must be matched by investments in local capacity; that transparency must be designed with privacy, equity, and historical bias in mind; and that infrastructural choices about vendors, governance, and ownership will determine whether enhanced access strengthens a shared evidentiary commons or simply adds another layer to existing power structures. For public interest data science, these are not peripheral concerns but design constraints: the value of mining municipal archives will be measured not only by what becomes searchable, but by who gains practical leverage from the resulting infrastructures and under what terms.

\subsubsection*{Synthesis}
Mining municipal archives is not just a vivid example of public interest data science; it helps specify what the field must be. Theoretically, it grounds abstraction about enclosure, exemption, and erosion in a concrete sociomaterial setting where law, infrastructure, and interface design jointly determine whether ``public'' data is actually usable. It also operationalizes the PIT×PIC meta-clusters: archives as evidentiary commons (journalism/law), audit substrate (accounting/engineering), infrastructural memory (planning), and narrative interface (civic tech and critical data work).

Practically, the case points to design patterns: machine-readable standards for agendas and minutes; open, documented ETL pipelines from PDFs to structured stores; redaction and access-control regimes co-designed with affected communities; and governance models that keep indexing, search, and AI-augmented retrieval under public or community control rather than proprietary dashboards. It highlights capacity-building as a core task: regional shared services, partnerships with universities and libraries, and reusable open-source toolchains.

As an agenda for public interest data science, mining municipal archives suggests concrete priorities: develop methods for ``audit-ready'' civic data pipelines; study the distributive effects of increased archival legibility; prototype worker–journalist–planner collaborations around local records; and advance institutional reforms that treat archival infrastructure as critical democratic infrastructure.

\subsection*{Case 2: Infrastructuring worker observatories}
% Gig work labor are also platforms with scant data and egregious enshittification dynamics
% Asymmetric power relations and data collection punishes workers with limited appeal
% Efforts to organize gig workers are creating observatories and folk infrastructures to collect and analyze their own data
% Legislation requiring platforms make these disclosures for transparency and implement appeals
% Clawing data back from exemption and implementing oversight
The platformization of work has reconfigured the basic terms of the labor contract. In ride-hailing, delivery, care work, and an expanding array of gig sectors, the most consequential information about workers' lives is collected, interpreted, and acted upon almost entirely by companies. Apps track GPS traces, acceptance rates, completion times, customer ratings, and opaque ``quality'' scores that determine access to shifts, bonuses, and even continued employment. At the same time, platforms themselves are subject to what critics call ``enshittification'': a gradual degradation of conditions for users and workers as firms optimize for investor returns. In this environment, data becomes both the means and the medium of control. Workers are constantly monitored, sorted, and evaluated, yet have almost no visibility into the metrics that govern them and few avenues for appeal.

The political economy of gig work is thus defined by extreme asymmetries of information and power. Platforms invest heavily in sociomaterial infrastructures—data centers, telemetry pipelines, dashboards, legal teams—that convert individual acts of labor into real-time behavioral streams. These infrastructures are tightly coupled to contractual regimes that classify workers as independent contractors, limiting obligations for benefits, job security, or due process. Deactivations, shadow bans, or re-ranking are frequently triggered by automated systems whose logic is treated as proprietary. Appeals, when available, are routed through support channels that offer little transparency or recourse. The sociotechnical configuration ensures governance outcomes that are predictably favorable to the firm: what counts as performance, what constitutes a violation, and who can be quietly excluded are all decisions embedded in metrics that the platform alone controls.

Against this backdrop, ``worker observatories'' have begun to emerge. Drawing on long traditions of workers' inquiry, mutual aid, and public interest organizing, gig workers and their allies are building grassroots data infrastructures designed to monitor, document, and analyze platform behavior from the bottom up. Driver associations, delivery collectives, and labor NGOs are experimenting with tools that allow workers to log pay, trips, wait times, tips, algorithmic changes, and deactivations. These may take the form of shared spreadsheets, mobile apps, browser extensions, or data diaries. Technically modest compared to corporate systems, these folk infrastructures perform a crucial epistemic function: they make it possible to detect patterns of unfair treatment, identify wage theft, evaluate the impact of algorithm updates, and substantiate claims in disputes with platforms or regulators. In effect, workers are attempting to build their own observatories to watch the systems that watch them.

\subsubsection*{Pressures and values}
The gig work context crystallizes the three forces identified earlier—enclosure, exemption, and erosion—and suggests how public interest data science might respond through openness, oversight, and ownership. Seen through this lens, infrastructuring worker observatories is not just a tactical innovation. It is a structural response to enclosure (by insisting on worker access and portability), to exemption (by creating independent oversight capacities), and to erosion (by building collective ownership and memory around worker data).

\paragraph*{Enclosure vs. openness}
Gig platforms are archetypal enclosures of labor data. Every ride, delivery, and rating is meticulously captured, but access to that information is tightly controlled. Workers see only thin slices of their own history; they rarely have access to the full distributions of pay, assignment patterns, or rating dynamics, and almost never to data about how others are treated. The platform's internal logs, model features, and experiments are walled off as trade secrets. Public interest data science counters this by reframing worker-generated data as a shared resource: openness here means ensuring that workers can access, export, and collectively analyze the information they produce, rather than letting it be locked away as proprietary capital.

\paragraph*{Exemption vs. oversight}
Exemption operates when platforms invoke intellectual property, terms of service, or jurisdictional ambiguity to avoid scrutiny. Algorithmic management systems that effectively make employment decisions are treated as outside the reach of labor law or administrative oversight because workers are nominally ``independent contractors'' and the code is proprietary. Worker observatories begin to claw back oversight by constructing independent vantage points. When hundreds of drivers log deactivations, sudden pay changes, or shifts in dispatch behavior, they create a counter-record that can be used to challenge corporate narratives, inform regulators, and support litigation. Oversight, in this frame, is not only a matter of state inspection but of building collective observability into systems that were designed to be one-way mirrors.

\paragraph*{Erosion vs. ownership}
Erosion appears in the precarity of worker knowledge and memory. Individual workers may notice that conditions are worsening—fewer high-paying jobs, more arbitrary penalties—but without records, those perceptions can be dismissed as anecdotal or idiosyncratic. Over time, turnover and churn erode institutional memory; lessons learned in one wave of organizing may be lost by the next. Ownership, as a public interest counter-force, involves building durable, worker-governed data infrastructures that outlast individual contracts and campaigns. A worker observatory that maintains longitudinal datasets, documented methodologies, and shared governance structures turns ephemeral experiences into a persistent knowledge commons. It embodies the idea that those who generate data should have a say in how it is stored, interpreted, and leveraged.

\subsubsection*{Meta-cluster mappings}
Worker observatories also map directly onto the four meta-clusters that emerged from the PIT×PIC grid search—themes that aggregate public interest lineages and technology conversations. In all four clusters, worker observatories demonstrate how public interest data science can operationalize the insights of journalism, law, accounting, engineering, and planning in a new domain: treating workers not merely as data subjects, but as co-authors of evidentiary infrastructures that hold platforms to account.

\paragraph*{Evidence, watchdogs, and rights}
At their core, worker observatories are evidence-making institutions. They document pay, conditions, and algorithmic behavior in ways that can support rights claims. Like investigative journalism, they collect and corroborate information about powerful actors that would otherwise remain hidden. Like public interest law, they generate the factual record needed for impact litigation, complaints to regulators, and collective bargaining. In doing so, they align squarely with the evidence-and-rights cluster: data is treated as a tool for contesting power, not just optimizing workflows.

\paragraph*{Metrics, audits, and institutional accountability}
Observatories also function as independent auditors. They compare ``promised'' earnings against actual incomes after expenses, track changes in platform fee structures, and examine how algorithmic tweaks redistribute risk and reward. This mirrors public interest accounting and engineering traditions that use metrics to test whether institutions are meeting their obligations. A well-run worker observatory is, in effect, a continuous audit of a platform's claims about flexibility, earnings, and safety—an instantiation of accountability metrics designed from the worker's vantage point.

\paragraph*{Infrastructures, space, and socio-technical safety}
Gig work is deeply spatial and infrastructural: drivers navigate city streets, couriers manage hazardous conditions, and algorithmic dispatch shapes exposure to risk. Worker observatories can incorporate spatial and temporal analysis—mapping where waiting times are longest, where accidents cluster, or how algorithmic changes shift traffic across neighborhoods. This situates them within the infrastructures-and-safety cluster: they help reveal how data-driven labor systems redistribute environmental and bodily risk, and can inform interventions by planners, regulators, and unions concerned with occupational health and urban justice.

\paragraph*{Interfaces, narratives, and data power}
Finally, worker observatories are about interfaces and narratives. The tools workers use—apps, dashboards, report generators—shape how they understand their own conditions and what levers they see as available. Stories built from observatory data can recalibrate public narratives around gig work, challenging marketing claims about entrepreneurship and flexibility with empirically grounded accounts of instability and control. This directly engages the interfaces-and-narratives cluster: observatories design touchpoints where workers, policymakers, and the broader public can encounter data about platform labor in ways that are intelligible, contestable, and tied to action.

\subsubsection*{Anticipated critiques}
Any serious proposal for worker data observatories has to contend with at least three persuasive lines of critique: questions about methodological rigor, concerns about safety and governance, and doubts about whether such infrastructures can meaningfully rebalance power rather than simply adding more work to already precarious jobs.

The first line of critique is methodological. Critics worry that worker-generated datasets will be noisy, partial, and systematically biased. Participation will almost certainly be uneven: those who are angriest or most engaged may be overrepresented; those in the most precarious positions may be least able to contribute. Data logged on phones or in shared spreadsheets can be incomplete, inconsistent, or misclassified. From a traditional data science perspective, these observatories may look like ``folk statistics''—difficult to validate, hard to generalize from, and vulnerable to both honest error and motivated reporting. If platforms and regulators can dismiss the resulting evidence as anecdotal or unscientific, skeptics ask, are we investing scarce organizing energy into tools that will have limited credibility in formal disputes? This critique does not deny workers' experience; it raises the bar for how observatories document methods, handle uncertainty, and connect qualitative testimony with quantitative aggregates in ways that withstand adversarial scrutiny.

A second, equally serious critique focuses on safety, privacy, and governance. Building a worker observatory means collecting detailed, time-stamped information about people's locations, earnings, and interactions with platforms—exactly the kind of data that, if mishandled, could expose contributors to retaliation or unwanted surveillance. Even well-intentioned organizers may lack the security expertise to protect sensitive logs; data leaks could endanger individuals or reveal organizing tactics. There is also the risk that observatories themselves become new sites of vulnerability or exploitation: third-party actors might seek to monetize the data, platforms might subpoena records, or factional disputes within worker groups might weaponize partial datasets. Skeptics argue that in already risky environments, asking workers to generate and share even more data about themselves may be ethically fraught unless governance structures, consent practices, and technical safeguards are robust and enforceable. On this view, the central question is not whether observatories are desirable in principle, but whether we are prepared to build them with the same level of care we expect from institutions handling medical or financial records.

The third line of critique is political and strategic: will worker observatories actually shift power, or will they become another layer of unpaid labor that platforms learn to route around? Critics note that many gig workers are already juggling multiple apps, unstable schedules, and high stress. Logging data, learning dashboards, and participating in analysis sessions may be feasible for a committed minority but unrealistic at scale. There is a danger that observatories primarily serve researchers, NGOs, or funders hungry for novel datasets, while the everyday conditions of work change little. Platforms may adapt by making their systems more opaque, adjusting terms of service to criminalize data collection, or co-opting the language of ``transparency'' with their own curated dashboards. Without clear pathways from observatory findings to bargaining power, regulatory enforcement, or public pressure, skeptics worry that these projects risk becoming technically impressive but politically marginal.

Taken together, these critiques do not argue against the idea of worker observatories; they specify what it would take for them to matter. They underscore that methodological rigor is not optional if the goal is to challenge powerful institutions; that safety and governance must be treated as core design problems, not afterthoughts; and that data infrastructures must be tightly coupled to organizing strategies, legal campaigns, and policy venues where evidence can be turned into leverage. For public interest data science, the challenge is to treat these concerns as design constraints: building observatories that are methodologically defensible, ethically governed, and strategically embedded in broader struggles over the future of platform work.

\subsubsection*{Synthesis}
Infrastructuring worker observatories shows what public interest data science looks like when the ``black box'' is an employer. Gig platforms enclose rich behavioral data about workers—routes, ratings, acceptance rates—while exempting their algorithmic management systems from meaningful oversight. Erosion appears as churn and amnesia: individual workers sense conditions worsening but lack durable records to prove it. Worker observatories respond by asserting openness (workers reclaiming access to their own data), oversight (independent monitoring of platform behavior), and ownership (collective stewardship of longitudinal worker-generated datasets).

Read through the PIT×PIC meta-clusters, the case is instructive. As evidence and rights, observatories produce records that underpin legal claims, organizing, and advocacy. As metrics and audits, they function as independent payroll and fairness audits, testing platforms' claims about earnings and flexibility. As infrastructures and safety, they surface spatial and temporal patterns of risk, from unsafe routing to extreme shifts. As interfaces and narratives, their tools and reports help reframe public stories about gig work from ``flexible opportunity'' to managed precarity.

Anticipated critiques—about data quality, privacy/security, and real leverage—sharpen the research agenda. Theoretically, they foreground worker-led sensing as a site for revisiting method, bias, and epistemic authority. Practically, they demand secure, usable, open-source toolchains and governance models co-designed with workers. Agenda-setting, they call for coupling observatories to regulatory reforms and labor institutions so that worker data becomes a durable counterweight to platform power.

\subsection*{Case 3: Anthropocene data trusts}
% Medium and long-term consequences of even best-case climate change imply massive disruptions that will also affect archives
% Current administration polices to shut down access to government datasets
% Obligations to future generations to preserve data that they can join with their own
% Building resilient archives and durable data infrastructures because data is both material and political
% Longtermist launching archives into space vs. commons-based ensuring archives remain open
% Clawing data back from erosion and implementing ownership
% Anthropocene data trusts thus exemplify public interest data science in its fullest sense. They treat data not simply as an input for models but as a long-term civic resource that requires institutions, infrastructures, and ethical commitments to endure. 
% Open data alone cannot counteract structural enclosure.
% Data is material, infrastructural, and shaped by power.
% Need for alternative, durable data infrastructures.
% Redundant archives.
% Community data trusts.
% Transparency registries.
% Civic data governance.

The Anthropocene is often framed in terms of energy transitions, infrastructure resilience, and political upheaval. Less visible, but equally consequential, is its challenge to how societies remember. Even under relatively optimistic climate scenarios, medium- and long-term impacts—sea-level rise, intensifying storms, wildfires, heat waves, and population displacement—will disrupt the ecological, social, and technical systems that sustain our collective ability to collect and archive records. Data centers depend on overtaxed power grids and vulnerable water supplies. Servers sit in floodplains and fire zones. Archives are housed in buildings never designed for intensifying climate extremes. At the same time, political decisions such as dismantling environmental data portals or defunding statistical agencies compound the risks of data loss.

The question is no longer only how we record temperatures, emissions, and land-use changes, but how we ensure that those records remain available and interpretable for decades and centuries to come. Future generations will need today's data not just to study the past, but to adapt, to allocate resources, and to press claims for responsibility and redress. If climate models cannot be validated because key observational records have vanished, or if evidence of environmental injustice has been quietly removed from public servers, the capacity for adaptation and accountability will erode just when it is most needed.

The concept of Anthropocene data trusts emerges at this intersection of climate disruption, data governance, and intergenerational justice. These proposed institutions treat data as both material and political: something that must survive physical shocks and institutional churn, and something that must be governed in ways that recognize obligations to people not yet born. Where traditional archives often center national heritage or institutional memory, Anthropocene data trusts foreground the continuity of ecological and social knowledge across deep time. They ask what it would mean to build archives capable of withstanding infrastructure failures, regime changes, commercial collapses, and natural disasters, while remaining accessible and meaningful to future publics facing different risks with different tools.

\subsubsection*{Pressures and values}
The climate crisis makes the triplet of enclosure, exemption, and erosion especially vivid and clarifies the countervailing commitments of openness, oversight, and ownership. Commitments to openness to resist enclosure, oversight to challenge exemption, and stewardship-oriented ownership to prevent erosion define the normative core of Anthropocene data trusts as a public interest data science project.

\paragraph{Enclosure vs. openness}
Climate-relevant data is increasingly enclosed. Satellite imagery is licensed to the highest bidder; proprietary models and sensor networks capture critical environmental signals that are then sold as subscription services; historical data once freely available via government APIs is restricted or paywalled. Even when datasets remain technically available, restrictive licenses and fragile interfaces limit meaningful reuse. Anthropocene data trusts push back by treating climate, environmental, and social records as non-excludable public goods. Openness here is not just about releasing CSVs; it is about ensuring that core evidentiary datasets like temperature records, emissions inventories, hazard maps, displacement statistics are maintained as part of a global commons, with open formats, clear documentation, and guaranteed access.

\paragraph{Exemption vs. oversight}
Exemption appears when powerful actors like states, corporations, even large NGOs claim the right to withhold or manipulate data on grounds of security, competitiveness, or liability. Emissions inventories may be incomplete, disaster data selectively reported, or model assumptions opaque. In the Anthropocene, such exemptions have generational consequences: decisions made today about what to disclose or suppress shape what future communities can know about who did what, where, and with what impact. Anthropocene data trusts seek to institutionalize oversight by creating independent custodial bodies that can demand, verify, and preserve critical datasets. They can serve as repositories for whistleblower disclosures, parallel archives for at-risk government records, and watchdogs that monitor when data disappears or degrades.

\paragraph{Erosion vs. ownership}
Erosion is perhaps the most insidious threat. It is the slow fading of bit rot, broken links, format obsolescence, budget cuts, and institutional neglect. A climate portal goes offline when a contract lapses; a regional database is lost during a reorganization; tapes and drives degrade unnoticed in an aging facility. Anthropocene data trusts reimagine ownership as stewardship: distributed, redundant, and intergenerational. They invest in geographically dispersed mirrors, migration plans for changing formats, governance mechanisms that outlast any single funding cycle, and legal structures that recognize obligations to future beneficiaries. Ownership in this sense is not about exclusive control; it is about shared responsibility to keep the record alive.

\subsubsection*{Meta-cluster mappings}
Anthropocene data trusts are not conjured from scratch. They synthesize lineages from journalism, law, accounting, engineering, and planning, and map naturally onto the four meta-clusters identified from the PIT×PIC grid: evidence and rights; metrics and audits; infrastructures and safety; interfaces and narratives. Across these clusters, Anthropocene data trusts look less like exotic data vaults in lunar caves and more like an evolution of familiar public institutions like libraries, archives, regulatory agencies, and newsrooms retooled for planetary-scale temporal and infrastructural challenges.

\paragraph*{Evidence, watchdogs, and rights}
At their heart, these trusts are about preserving evidence. Like investigative journalists, they safeguard records that may support future claims about responsibility and harm: emissions data, regulatory correspondence, community surveys, disaster responses. Like public interest law, they recognize that rights to clean air, water, and a stable climate are hollow without enduring proof of who polluted what, where, and when. Anthropocene data trusts build the evidentiary commons that future watchdogs, litigators, and truth commissions will need.

\paragraph*{Metrics, audits, and institutional accountability}
Long-term environmental data is the backbone of many accountability mechanisms: emissions caps, biodiversity targets, adaptation benchmarks. Public interest accounting and engineering traditions remind us that metrics matter only if they are auditable. An Anthropocene data trust can maintain reference datasets and documentation that allow future analysts to verify whether reported reductions were real, whether offsets were backed by actual sequestration, and whether adaptation funds reached vulnerable communities. In this sense, they are institutionalizing the metrics-and-audits cluster on a planetary scale.

\paragraph*{Infrastructures, space, and socio-technical safety}
Climate change is fundamentally about infrastructures and space: sustaining human and natural systems in the face of increasingly irreversible climate shocks. Anthropocene data trusts must therefore operate as infrastructural actors. They rely on resilient physical facilities, network connectivity, and energy sources, and they curate spatial datasets like land-use histories, hazard maps, migration flows that shape where risk accumulates and which communities are visible in planning. They align with the infrastructures-and-safety cluster by treating data archives themselves as critical infrastructure for climate adaptation and disaster governance.

\paragraph*{Interfaces, narratives, and data power}
Finally, archives are only as powerful as the stories they enable. Anthropocene data trusts must design interfaces and narrative strategies that make deep-time data usable and meaningful: tools that allow future planners, journalists, and citizens to trace the lineage of a decision, reconstruct the consequences of disasters, or understand patterns of displacement. They must grapple with representation: whose experiences are recorded, which languages are included, what metadata frames events. This situates them squarely in the interfaces-and-narratives cluster: they shape how past data is encountered, interpreted, and mobilized in future debates about justice and responsibility.

\subsubsection*{Anticipated critiques}
A project as ambitious as Anthropocene data trusts will inevitably attract serious and healthy skepticism. Rather than treating these objections as noise, it is useful to see them as three intertwined lines of critique that any credible proposal must confront: feasibility, politics, and power.

The first concern is feasibility in the broadest sense: can we realistically build and sustain such institutions over the time scales they imagine? Critics point to the fragility of existing archives, the chronic underfunding of public infrastructure, and the sheer technical complexity of managing petabytes of heterogeneous data over decades and centuries. If universities struggle to maintain their own repositories and governments routinely allow important portals to languish, why should we believe that a new class of ``Anthropocene data trusts'' will fare better? This line of critique cautions against wishful thinking and imagining elegant governance structures that are never fully resourced or politically defended. A public interest response has to take this seriously, grounding proposals in incremental, federated models that build on institutions we already have like libraries, regulatory agencies, and community archives rather than assuming the emergence of a well-intentioned global steward.

The second critique centers on politics and representation: who decides what is preserved, on what terms, and with whose knowledge? Here, critics worry that Anthropocene data trusts could quietly reproduce existing hierarchies: archiving the data of powerful actors and well-resourced regions while marginalizing Indigenous knowledge, local records, or inconvenient evidence of injustice. Decisions about which datasets are ``canonical,'' which vocabularies define the archive, and which communities are framed as data subjects versus co-governors are deeply political. Without careful design, a data trust meant to serve intergenerational justice could entrench a narrow, technocratic view of what counts as climate-relevant knowledge. This critique pushes public interest data science to foreground participatory governance, to embed principles from data justice and decolonial scholarship, and to treat the trust not as a neutral vault but as a contested space where questions of inclusion, consent, and epistemic authority are continually negotiated and documented.

A third, related line of critique focuses on power and cooptation. Even if the technical and participatory hurdles can be addressed, skeptics ask whether Anthropocene data trusts might be captured by the very actors whose behavior they are meant to make legible. Corporations might support a trust as a reputational shield while steering it away from datasets that would enable meaningful accountability. States could use the language of intergenerational stewardship to centralize control over environmental data and restrict access under the guise of security. There is also a risk that ``trust'' language masks sharp conflicts over intellectual property and commercial exploitation, especially in a post-AI landscape where training data has become a lucrative asset. In this view, the danger is not only that the archives fail, but that they succeed on terms that turn a planetary memory project into another opportunity to enclose the commons.

Taken together, these critiques do not refute the idea of Anthropocene data trusts; they specify the conditions under which such institutions might be worth building. They remind us that durability without funding is fantasy, preservation without pluralism is injustice, and trust without countervailing power is na\"{i}ve. A public interest agenda must therefore treat feasibility, politics, and power as design constraints rather than afterthoughts using them to sharpen which data is prioritized, which partnerships are pursued, and which safeguards are non-negotiable. Anthropocene data trusts should reflexively serve future publics rather than simply inheriting the blind spots and asymmetries of the present.

\subsubsection*{Summary}
Anthropocene data trusts make clear that public interest data science must grapple with climate risk as a problem of memory and governance, not just modeling. The motivation is stark: medium- and long-term climate impacts, political interference, and infrastructural fragility threaten the ecological, social, and regulatory records future generations will need for adaptation and accountability. Enclosure (proprietary climate and emissions data), exemption (opaque reporting, suppressed datasets), and erosion (bit rot, link rot, institutional churn) all push against values of openness, oversight, and stewardship-oriented ownership.

Mapped onto the PIT×PIC meta-clusters, Anthropocene data trusts are: evidence and rights infrastructures preserving proof of responsibility and harm; metrics and audits backbones for verifying climate targets and financial claims; infrastructures and safety for informational resilience alongside physical adaptation; and interfaces and narratives that mediate how future publics encounter the past. Anticipated critiques about feasibility, representation, and co-optation sharpen the agenda: durability without funding is fantasy; preservation without pluralism entrenches injustice; ``trusts'' without safeguards risk becoming new sites of enclosure.

Theoretically, this case demands a long-horizon, intergenerational lens on data commons and justice. Practically, it implies federated archival architectures, mirrored climate and regulatory datasets, and governance charters co-designed with affected communities. As agenda-setting, it calls for treating climate-relevant data infrastructures as critical public utilities, with dedicated regulation, funding, and professional roles anchored in public interest data science.

% \section*{Public interest data futures}
% Public interest data science offers methodological and institutional tools for transparency.
% Ensures accountability and justice amid expanding technological power.
% Calls for resilient data commons and regulatory frameworks that safeguard democratic oversight.

% Advancing public interest data futures requires both top-down and bottom-up strategies that treat data access and stewardship as core democratic obligations rather than peripheral technical concerns. At the level of law and regulation, this means updating legislation to protect and expand data access rights by drawing explicitly on established regulatory models. Just as the SEC mandates financial disclosure and auditing to guard against fraud and systemic risk, public interest data regimes should require transparent reporting and independent evaluation of algorithmic systems. Analogous to the FCC's role in ensuring communications fairness and nondiscrimination, regulators should enforce baseline standards for platform transparency, access to public-relevant data, and the ability of researchers and civil society to scrutinize information ecosystems. The FDA's emphasis on safety and efficacy in public-impact systems offers a template for pre-deployment testing and post-market surveillance of high-stakes AI applications, while the EPA's mandate to monitor and mitigate externalities underscores the need to track and remediate social and environmental harms produced by data-driven infrastructures. These top-down measures institutionalize openness, oversight, and ownership as public obligations rather than corporate options.

% Legal change alone, however, is insufficient without corresponding shifts in professional norms and funding priorities. Public interest data futures depend on valuing archival work, data stewardship, and accessibility as central contributions rather than low-status maintenance tasks. Universities, funding agencies, and professional associations must reward the creation of durable data repositories, high-quality documentation, and open infrastructures alongside methodological innovation. This will also require creating professional pathways to support practitioners moving across institutional settings like government, academia, civil society, and journalism without abandoning public interest commitments. Developing credentials, fellowships, and career ladders explicitly oriented toward public interest data science can help stabilize this work, making it possible for practitioners to maintain continuity across projects and sectors. Only by aligning regulation with professional practice of treating data as a long-term civic resource rather than disposable exhaust can public interest data science build the stable foundations required for democratic oversight.

% Bottom-up priorities are equally essential. They depend on treating communities not as data subjects but as co-governors of data infrastructures. One starting point is to adapt models of decentralization and participatory governance from cooperatives, unions, and community trusts. In these traditions, members collectively decide how resources are managed, how benefits are distributed, and how risks are mitigated. Applied to data, this suggests structures such as community data trusts, worker-led observatories, and municipal data cooperatives in which decisions about collection, access, and use are made through deliberative, representative processes rather than solely by corporate or state actors.

% Equally important is the cultivation of civic and rhetorical norms that prioritize the use of data in public deliberations and debates, not simply as technocratic evidence but as shared infrastructure worth sustaining. When residents, journalists, organizers, and policymakers routinely invoke public data to argue about budgets, zoning, policing, health, or climate, they generate constituencies for preserving and improving those data resources. The act of using data in public argument becomes, itself, a justification for resisting enclosure and investing in stewardship.

% Public interest data futures require thinking about data not as a static asset but as an evolving relationship among institutions, communities, and infrastructures. Top-down reforms without grassroots capacity risk producing formal rights that few can exercise; bottom-up innovation without legal guarantees risks being overwhelmed by enclosure, exemption, and erosion elsewhere in the ecosystem. The challenge is to align regulatory architectures with participatory governance models: to design laws that anticipate and support cases like municipal archives, worker observatories, and data trusts rather than treating them as peripheral experiments.

% Public interest data science is less a discrete specialty than a re-alignment that cuts across disciplines and sectors. It calls on data scientists, computer scientists, statisticians, and information professionals to collaborate with journalists, lawyers, organizers, planners, and archivists in sustaining data as a civic resource. It asks funders to treat maintenance and stewardship as intellectually substantive work, and it invites educators to train practitioners who can move fluidly between contexts. Public interest data futures will be built incrementally through new statutes and standards, new archives and observatories, new professional pathways and public narratives: reclaiming the conditions under which data-intensive systems can be subject to similar values, rules, and norms governing other risky and high-impact infrastructures.


\section*{Discussion}
If data-driven systems are now core infrastructures of public life, then we must be able to see, question, and govern them. Today, that capacity is being undermined by three interlocking pressures. Enclosure fences off data that once functioned as a commons, from social media archives to environmental and administrative records. Exemption allows dominant platforms and AI developers to operate outside meaningful regulatory scrutiny, invoking trade secrecy, privacy, and outdated safe-harbor doctrines to avoid disclosure and accountability. Erosion quietly degrades the institutional and technical infrastructures like APIs, public datasets, archives, practices of documentation that once made long-term observation and reproducible analysis possible.

Public interest data science is a response that takes these pressures as design constraints rather than background noise. It insists on building countervailing commitments of openness, oversight, and ownership. These are not abstract virtues but are grounded in the public-interest traditions of other technical disciplines and professions like journalism, law, accounting, engineering, and planning, and informed by adjacent conversations in civic tech, data-for-good, and critical data studies. Journalism's watchdog ethos, public interest law's concern with structural inequality, accounting's auditing standards, engineering's safety culture, and planning's advocacy practices all show how professions have historically turned technical expertise toward the public good. Civic tech and tech-for-good communities add tool-building and experimentation; critical data studies and data justice remind us to foreground power, inequality, and epistemic harm.

Openness counters enclosure by treating society's ability to observe sociotechnical systems as a necessity, not a luxury. In practice, this means investing in redundant public archives, machine-readable standards, transparent pipelines, and community-governed data trusts so that key records do not vanish when corporate strategies or political winds change. Oversight responds to exemption by developing methods, standards, and institutions capable of evaluating high-impact algorithms and data infrastructures, drawing analogies to financial audits, product safety testing, and environmental impact assessment. Ownership, conceived as stewardship rather than exclusive control, addresses erosion by distributing responsibility for long-term data maintenance across public institutions, civil society, and affected communities, with explicit obligations to workers and future generations.

Our RAG-augmented, grid-search analysis of public interest traditions and public interest technology conversations surfaced four meta-clusters that organize this agenda: Evidence and Rights, Metrics and Audits, Infrastructures and Safety, and Interfaces and Narratives. Together, they mark out a practical program. Evidence and Rights focuses on building data infrastructures that support claims-making, legal redress, and democratic oversight. Metrics and Audits concerns the design of indicators, reporting regimes, and audit methods that make institutional performance legible and contestable. Infrastructures and Safety treats data and algorithmic systems themselves as critical infrastructures that must be resilient, equitable, and governable. Interfaces and Narratives addresses how people encounter, interpret, and act on data—through portals, dashboards, reports, and stories.

\subsection*{Speculative cases}
The three case studies give this framework concrete texture. Mining municipal archives shows how legally public records like council packets, budgets, permits, enforcement logs are effectively enclosed by poor standards, fragmented systems, and brittle interfaces. Public interest data science here means designing ``audit-ready'' civic data pipelines, AI-augmented retrieval tools under public or community control, and archival infrastructures that make local government genuinely observable over time. Infrastructuring worker observatories situates public interest data science in the world of platform labor, where companies enclose rich behavioral data and exempt algorithmic management from scrutiny. Worker-led observatories reclaim openness (access to one's own data), build oversight (independent monitoring of pay, risk, and fairness), and experiment with collective ownership of longitudinal worker-generated datasets as a counterweight to platform power. Anthropocene data trusts extend the argument into intergenerational time: climate disruption and political volatility threaten to erase ecological, regulatory, and social records exactly when future publics will need them most. Here, the task is to build federated, commons-oriented archival institutions that preserve climate data as evidence for future generations, as accountability infrastructure, and as a shared resource for adaptation.

Across these sites, a common set of critiques recurs, and any credible agenda must take them seriously. A first line of critique asks about feasibility: can under-resourced municipalities, precarious workers, or overstretched archives realistically sustain the infrastructures we imagine? A second line interrogates politics and representation: who decides what is collected and preserved, whose knowledge is incorporated, and who is rendered newly visible or vulnerable when data becomes more searchable and linkable? A third line focuses on power and co-optation: might archives, observatories, or data trusts themselves be captured by the actors they are meant to hold accountable, or become new channels for commercial enclosure and ``transparency theater''?

\subsection*{Implications for theory, practice, and institutions}
In this sense, anticipated critiques do not weaken the case for public interest data science; they sharpen its contours. They remind us that durability without funding is fantasy; that preservation without pluralism can reproduce injustice; that ``openness'' without attention to privacy and harm can deepen surveillance; and that ``trusts'' or observatories without strong governance can become new instruments of control. Public interest data science must therefore be realist, critical, and pragmatic: designing institutions that are resource-aware, participatory, and resistant to capture, and building methods that can withstand adversarial scrutiny.

What, then, does this imply as a closing agenda for theory, practice, and institution-building? For theory, the cases and meta-clusters call for a more explicit synthesis of sociotechnical, legal, and political-economic analysis. We need frameworks that treat data infrastructures simultaneously as technical systems, as legal-regulatory objects, and as contested commons, with attention to how enclosure, exemption, and erosion operate over time and how openness, oversight, and ownership can be concretely institutionalized. Public interest data science should be a site where questions of democracy, labor, and intergenerational justice are treated as first-order concerns, not afterthoughts to optimization.

For practice, the agenda is to prototype and evaluate concrete patterns: civic RAG systems that mine public archives without privatizing them; worker-led sensing and logging tools with robust security and governance; federated climate and regulatory archives with clear charters, redundancy plans, and community representation. It includes developing audit methodologies that work with partial, externally observable data; codifying standards for ``public interest–ready'' data pipelines; and building open-source, reusable toolchains that lower barriers for municipalities, unions, newsrooms, and NGOs.

For institutions, the agenda is to align regulation, professional norms, and grassroots capacity. Regulatory reforms should adapt lessons from the SEC, FCC, FDA, and EPA to data and AI—creating enforceable obligations for transparency, documentation, and independent evaluation of high-impact systems. Universities, agencies, and funders should treat data stewardship, audit work, and community partnership as central scholarly and professional contributions. New career paths, fellowships, and credentials should enable public interest data scientists to move fluidly across academia, government, civil society, and journalism. And civic cultures should normalize the expectation that data—about budgets, platforms, and emissions—is available, interrogable, and used in public argument.

\section*{Conclusion}
Building such a field will not be simple but the stakes are clear. Data-driven systems already shape the conditions of work, the texture of public life, and the horizons available to future generations. A public interest data science, anchored in openness, oversight, and stewardship, offers one path to ensuring that those systems are answerable not only to shareholders and clients, but to the public, to workers, and to future generations who will live with the consequences of our decisions.

\begin{acks}
Thanks for all the fish.
\end{acks}

%%
%% The next two lines define the bibliography style to be used, and
%% the bibliography file.
\bibliographystyle{ACM-Reference-Format}
\bibliography{bibliography}

\end{document}
\endinput
%%
%% End of file `sample-acmsmall-submission.tex'.
